\documentclass[10pt,a4paper]{amsart}
\usepackage[margin=19mm]{geometry}
\usepackage{amsmath,amssymb,mathtools}
\usepackage[scr=rsfs,cal=euler]{mathalfa}
\usepackage{xfrac}
\usepackage[unicode,hidelinks,bookmarks=false]{hyperref}
\newcommand{\ie}{i.e.\ }
\newcommand{\cf}{cf.\ }
\newcommand{\resp}{respectively\ }
\newcommand{\Subsection}{\subsection}
\providecommand{\nobreakdash}{\nobreak\hbox{-}}
\setlength{\emergencystretch}{2em}
\multlinegap0pt
\usepackage{mathtools}
\usepackage{cancel}
\usepackage{amssymb}
\usepackage{tikz}
\usepackage{enumerate}
\usepackage{float}
\usepackage{dsfont}
\usepackage{xcolor}
\newtheorem{theorem}{Theorem}
\newcommand {\R}{\mathbb R}
\title{\LARGE \bf An application of the mean motion problem to time-optimal control}
\author{Omri Dalin, Alexander Ovseevich, Michael Margaliot}
\date{Source: arXiv:2502.13523v1 (CC BY 4.0)}
\begin{document}
\maketitle

\noindent\emph{Source and licence.} \LARGE \bf An application of the mean motion problem to time-optimal control. Version arXiv:2502.13523v1 (CC BY 4.0), distributed by arXiv under the Creative Commons Attribution 4.0 International licence, http://creativecommons.org/licenses/by/4.0/, as recorded in the arXiv licence metadata for this exact version and shown on the abstract page https://arxiv.org/abs/2502.13523v1. Full licence notice: \texttt{licenses/arxiv-2502.13523-CC-BY-4.0.txt}.\par\smallskip

\noindent\textit{Editorial scope.} Counted unit: the article's theorem imaginary (source0.tex lines 1108--1110). The article's complete original proof of it is delivered with it (source0.tex lines 1111--1114, one \texttt{proof} environment of 4 lines). No mathematics is written, repaired or re-derived: the statement and the proof are the article's own lines, in the article's own notation, and no retained line is substituted or re-flowed.  Retained uncounted as the source-local context the proof needs: lines 1059--1061.  Nothing was omitted from any retained span, so the package carries no omission record.  No retained line contains a citation, so the wrapper carries no bibliography and no bibliography entry.  No retained line contains a figure environment, an image-inclusion command or drawing code, so no image is delivered and the package declares no asset dependency.  Editorial formatting only, in addition: an AMS \texttt{amsart} preamble that restates the article's own declarations and macros used by the retained lines, namely the environments \texttt{theorem} on separate counters, together with the standard packages those lines need.  The retained environments print in this document as Theorem 1 in order of appearance, whereas the article prints the same items with the numbering of its own full document.  The complete native source of the article as distributed in the arXiv e-print for this exact version is bundled unchanged as \texttt{sources/173523/source0.tex}.
\begin{theorem}\label{lower}
The number $N(T)$ of zeroes of the switching function in the interval $[0,T]$ has the lower bound $N(T)\geq CT$, where $C$ is a positive constant, as $T\to\infty$.
\end{theorem}

 \begin{theorem}\label{imaginary}
Suppose that the linear system is controllable, and all the eigenvalues of the matrix $A$ with the maximal real part  are not real. Then, for a generic $ p\in\R^n$ the number of roots of the function $t\mapsto p^\top e^{At}b$ in the interval $[0,T]$ of time is greater that $cT$. Here $c>0$ is a strictly positive constant, and $T>0$ is sufficiently large.
\end{theorem}

\begin{proof}
 This result gives a linear lower bound for the number of switching points for a control system $\dot x=-Ax+bu$, with the governing matrix $A$, since the switching function for the system  is $p^\top e^{At}b$.

This can be proved by a straightforward generalization of the proof of Theorem \ref{lower}. Indeed, consider the $L_1(0,T)$-norm of the function $g(t):=e^{-\nu t}f(t)$, where $\nu$ is the maximal real part of eigenvalues of $A$, and $f$ is the switching function. Then, if the ``dual vector'' $p$ has a non-zero component in the root space of $A^\top$, corresponding to an eigenvalue with the real part $\nu$,   the leading term in  the $L_1(0,T)$-norm of $g$ looks like $\int_0^T|\sum a_k(t)e^{i\mu t}|dt$, where $\mu$ runs over the imaginary parts of the eigenvalues with the real part $\nu$, and $a_k$ is a polynomial in $t$. This expression is studied in the proof of Theorem \ref{lower} and is shown to be greater than $ cT^{d+1},\, c>0,$ where $d$ is the maximal degree of polynomials $a_k$. On the other hand, the $\int_0^T\sum a_k(t)e^{i\mu t}dt=O(T^d)$. These two inequalities, pointing in opposite direction, are incompatible. This  shows that there is a zero of $g(t)$ in any sufficiently long interval of time.  Therefore, the number of zeroes of $g(t)$ in $(0,T)$ grows linearly with $T$. Since the zeroes of $g$ are exactly the same as that of the switching function $f$, we get the required lower bound on the amount of switches. \end{proof}

\end{document}
